\section{The Labors of Hercules, Part 5}

The Augean Stables\footnote{\url{https://en.wikipedia.org/wiki/Augeas}}

Hercules does not go immediately into action or battle, following his insult from the gods: Hera's madness that causes him to murder his own offspring is the ``fall'' or primordial condition of sinful man who is an enemy to himself and his own kith/kin, let alone everyone else -- man as a ``wolf to other men''. Finding one's self in this condition requires ``penance'', but is really not the loss of anything, since it is illusory to begin with -- the hero doesn't find penance distressing. The path out of his dilemma is not ``out'', but ``deeper in and through''. This we have seen -- The Man awakes to his sinful or mad condition, \& rather than rage against the gods, he shoulders his karmic burden, such as it is, channeling his energies into carefully executing the appointed labors. So the Leftist solution of railing at the universe or shaking one's fist at the sky is a non-starter -- \emph{change begins with the Self}. But even further, with some change begun, one does not therefore jump up and ``clean out the Augean stables'', by reforming the Church, cleansing the State, defeating the forces of modernization and liberalism, etc., etc. It is much easier, when awakened, to immediately but covertly ``rage against the machine'' by marching off to clean out the Augean stables: let's reform the Church of Rome! Let's get rid of our traitors at home! Let us take counsel together, and purge the earth of the presence of all that is impure!

No, Hercules has worked up to this condition, \& is being given tasks to perform in a logical and precise sequence. He doesn't nominate or choose himself for any of them -- he is the man for the job, and even the gods recognize this, but the jobs are being chosen proximally by his enemies. He is submitting to government through Providence, and trusting in a higher power than himself. This, of course, means that within the Hercules-cycle, there will finally be no ``higher Power'' but Hercules, but that is for the final chapter, upon his exaltation as a constellation.

The Augean stables are a problem that everyone knew about, and no one had solved. We could say that this is ``the modern world''. What in the world do we do with all the refuse and offal from the ``madding crowd'' which inhabits the round globe, one and all having fallen prey to madness and imbecility? How will this impurity be cleansed?

Once again, Hercules does not come raging in headlong, like a Conan the Barbarian or Pol Pot, bent intently on ridding the earth of ``undesirables''. He uses the forces of Nature itself, altering their course, to bring together the natural remedy for such a large heap of trash and garbage: he seizes a river and changes the course, immediately washing the stables.

He does this, presumably, for gratis, as the labor will be determined a ``failure'', since he did not rely on his own two hands to do it. Hercules has been more cunning than those, however, who cunningly disqualify the trick. He has succored fellow man and beast, through the bounty of Nature, and provided a tangible good to all of those who are beneath him. He has stooped and entered the world of ``common man''. This is what makes of Hercules more than a hero and an adventurer, and begins to give him the title of benefactor or popular hero. He will enter the \emph{mythos} as a man like unto King Alfred, who brought permanent peace to England by sponsoring and converting a noble pagan, Guthrum, thus creating the Danelaw and bringing a new era of harmony, although he Alfred was redoubtable in battle nonetheless. The gods and test makers disqualify this test, but one can see that Hercules would have done it anyway. Why? \emph{Because it needed doing}. Often it happens that those who make the ``most progress'' esoterically are those who are least concerned about it, because detachment really means detachment. Have you ever heard of the old Puritan preacher who found peace in the thought that he would at least be a monument to God's justice in Hell?

\emph{It may be that the gulfs will wash us down...}\footnote{\url{http://www.poetryfoundation.org/poem/174659}}

Tennyson has Ulysses say this, looking out over his realm.

That is the spirit in which to embark upon dread tests, to balance out the focus and energy one brings to them (also requisite). One must play the game because the game matters. In esoteric matters, it really is ``how you play the game''. Intention and personal goals, rather than a lust for power, are relevant. Even personal salvation takes a back seat to the ``point'' of what is happening. A hero who has no time to ``stoop'' and assist someone is a hero who is worried too much about husbanding their strength. Spend it, and it will be given again. Whatever your hand finds to do, do it with all might. That is the path of Hercules.

His digression is actually part of the Zodiac pattern, which must in any case have 12, not 10 labors.\footnote{\url{http://www.thelemapedia.org/index.php/The\_Twelve\_Labors}}

An important point Cologero has made over and over is that esoteric studies is not hostile to exotericism. This, in itself, is more than likely a result of the action of God upon the human psyche, as expressed in Christianity. It is now revealed that a man not only can, but ought, to pursue esoteric studies in harmony with the \emph{Logos}: that is, if you desire enlightenment, you should not neglect the study of what ``sane and noble men'' of the past thought, how they comported themselves, and what they believed. It is permissible to ``prepare for enlightenment'' in this way -- by finding mentors (in the past if necessary). If done properly (in a healthy way), this does not ``add'' to psychic accretions or burdens, but functions as a means of restoring normal function in the psyche.

For Hercules, this is accomplished dramatically in his willingness to do an altruistic and noble deed, even though it prolongs his labors, formally. He stoops to conquer, and aids the disharmony of Nature, with sheer physical power and the elements at his disposal. In modern language, this is the much derided ``helping a little old lady across the street''. It is the Earps cleaning up Tombstone. It is Oliver Cromwell threatening to send English warships against the Turks. It is the hero standing against a playground bully on behalf of someone quite small. The hero, here, seems to forsake his task and disqualify himself -- he is ``distracted'' by the enormity of evil around him. In reality, he simply does what a real man would do, if he comes across that which he cannot abide. This comes later in the growth of the warrior-saint. It is not an immediate goal, but emerges organically as he progresses. He is, in reality, being swiftly drawn out of his esoteric path, by a deeper current, which is directly connected to the Divine. This current only seems to cost him in the short run. In actuality, it defines the essence of what he will become, although he is not yet purely and entirely ready. It is reflected frequently in popular culture\footnote{\url{https://www.youtube.com/watch?v=6aHkRn9UL2Q}}, whose subconscious is aware of much of these truths.

It is analogous to distracting the already distracted psyche by giving it a ``bone to chew on''. In esoteric practice, this is Ora et Labora\footnote{\url{https://en.wikipedia.org/wiki/Navigation\_Acts}}. This is the power that gave rise to the dominance of Christendom's arms, time and again, against much greater odds, when by all rights, Europe should have crumbled, collapsed, or shattered under the hammer blows of a much more ``this-worldly'' religion -- Islam during its height. The religion of silver, and of the crescent, was better organized, less complex dogmatically, more victorious in arms, and oriented towards secular power and monolithic outer form. Time and again, the waves of its seas lapped at the fragile edges of Christendom, and time and again, Christendom provided men who put a pause to their own progress to take up the sword in behalf of the altar. In secret studies, one one would meditate on images, icons, or sacred ``words'' or phrases given one by the Spirit, while one is learning to concentrate the power of the nous, first of all towards the body, then, into the psyche, and finally into the nous-itself, prior to re-divinization in the ``Self beyond the Self'', as the nous travels back to God. Robin Amis teaches that the double arrow, or the straight line and the circle, symbolizes the power of the nous that is ``in the world, but not of the world''. This might also be the double-headed eagle, so commonly seen in European royal symbolism.

Here, Hercules meditates on a relatively simple and mundane problem, but one that no one has been able to solve. He came, he saw, he conquered. It doesn't matter what the gods decree -- Hercules will finish, and the labor will count unofficially, if nothing else, in the popular mythology of the folk, who recognize the good deed for what it is. In any case, Hercules is now beginning to have ``power to spare''.

A poem beautifully expressing this is Thomas Moore's \emph{Prince's Day}\footnote{\url{https://www.youtube.com/watch?v=ogYjDfuf5S0}}.

The body and psyche, ``in chains'' and suffering, nevertheless expresses a strong yearning for the will of the Prince. Paradoxically, the suffering and yearning of the body/psyche provide a ``distraction from distraction'' -- the nous is turned farther inward, and the path rather than being deflected, is intensified. The ocean wave withdraws to gather its strength. I include the whole poem -- it is part of the Logos, the gift of the true West to men. When the book of Kells was created, something was added to the religion of Christ (I say ``added'', but rather revealed as latent). The East isolated images in their essential purity, but if you look at the pages of the book of Kells\footnote{\url{https://en.wikipedia.org/wiki/Book\_of\_Kells}}, you see knot-work, playful animals and creatures, and a kind of background dance and delight in Nature. The Celtic and Frankish saints who converted the West were immersed in the world, and yet, still, not of it.

\paragraph{The Prince's Day}

\begin{verse}
Though dark are our sorrows, today we'll forget them,\\
And smile through our tears, like a sunbeam in showers:\\
There never were hearts, if our rulers would let them,\\
More form'd to be grateful and blest than ours.\\
But just when the chain, Has ceased to pain,\\
And hope has enwreathed it round with flowers,\\
There comes a new link, Our spirits to sink ---\\
Oh! the joy that we taste, like the light of the poles,\\
Is a flash amid darkness, too brilliant to stay;\\
But, though 'twere the last little spark in our souls,\\
We must light it up now, on our Prince's Day.

Contempt on the minion who calls you disloyal!\\
Though fierce to your foe, to your friends you are true;\\
And the tribute most high to a head that is royal,\\
Is love from a heart that loves liberty too.\\
While cowards, who blight Your fame, your right,\\
Would shrink from the blaze of the battle array,\\
The Standard of Green In front would be seen ---\\
Oh, my life on your faith! were you summon'd this minute,\\
You'd cast every bitter remembrance away,\\
And show what the arm of old Erin has in it,\\
When roused by the foe, on her Prince's Day.

He loves the Green Isle, and his love is recorded\\
In hearts which have suffer'd too much to forget;\\
And hope shall be crown'd, and attachment rewarded,\\
And Erin's gay jubileee shine out yet.\\
The gem may be broke By many a stroke,\\
But nothing can cloud its native ray;\\
Each fragment will cast A light to the last ---\\
And thus, Erin, my country, though broken thou art,\\
There's lustre wiithin thee, that ne'er will decay;\\
A spirit which beams through each suffering part,\\
And now smiles at all pain on the Prince's Day.
\end{verse}



\flrightit{Posted on 2013-09-19 by Logres}